% year: תשפ"ד
% type: מבחן
% semester: א
% moed: מיוחד
% number: 2
% points: 16
% categories: continuity, func-limits
% source: exams/מבחנים/תשפד/מועד ג תשפד.pdf

\begin{question}
נגדיר פונקציה על ידי
\[ f(x)=\begin{cases}2^{\frac{\sin x}{x}} & x<0,\\ ax+b & 0\le x\le 1,\\ \frac{x^2+x-2}{\sqrt{x}-1} & x>1\end{cases} \]
לאילו ערכים של הפרמטרים $a,b$ הפונקציה רציפה בכל $\R$? נמקו.
\end{question}

\begin{hint}
בתוך כל תחום הפונקציה רציפה; נשאר לדרוש רציפות בנקודות התפר $x=0$ ו-$x=1$.
\end{hint}

\begin{hint}
עבור $x>1$: $x^2+x-2=(x-1)(x+2)$ ו-$x-1=(\sqrt x-1)(\sqrt x+1)$.
\end{hint}

\begin{solution}
\textbf{בתוך התחומים.} ב-$(-\infty,0)$ הפונקציה $\frac{\sin x}{x}$ רציפה (מנה של רציפות, מכנה שונה מ-$0$), ו-$2^t$ רציפה, ולכן ההרכבה רציפה. ב-$(0,1)$ זו פונקציה לינארית, רציפה. ב-$(1,\infty)$ זו מנה של פונקציות רציפות שהמכנה שלה $\sqrt x-1>0$, ולכן רציפה. לכן לכל $a,b$ הפונקציה רציפה בכל נקודה פרט אולי ל-$x=0$ ו-$x=1$.

\textbf{רציפות ב-$x=0$.} $f(0)=b$, ומימין $\lim_{x\to0^+}(ax+b)=b$. משמאל, מאחר ש-$\lim_{x\to0}\frac{\sin x}{x}=1$ ו-$2^t$ רציפה ב-$t=1$, לפי משפט הגבול של פונקציה מורכבת
\[ \lim_{x\to0^-}2^{\frac{\sin x}{x}}=2^1=2. \]
לכן $f$ רציפה ב-$0$ אם ורק אם $b=2$.

\textbf{רציפות ב-$x=1$.} $f(1)=a+b$ ו-$\lim_{x\to1^-}f(x)=a+b$. מימין, עבור $x>1$:
\[ \frac{x^2+x-2}{\sqrt x-1}=\frac{(x-1)(x+2)}{\sqrt x-1}=\frac{(\sqrt x-1)(\sqrt x+1)(x+2)}{\sqrt x-1}=(\sqrt x+1)(x+2), \]
ולכן $\lim_{x\to1^+}f(x)=(1+1)(1+2)=6$. לכן $f$ רציפה ב-$1$ אם ורק אם $a+b=6$.

\textbf{תשובה:} $f$ רציפה בכל $\R$ אם ורק אם $\boxed{b=2,\ a=4}$.
\end{solution}
